\section{特征值与特征向量}

\subsection{特征值问题}

\begin{BoxDefinition}[特征值与特征向量]
    对于$\vb{A}\in\R^{n\times n}$，若存在$\vb{v}\in\C^n$且$\vb{v}\neq\vb{0}$以及$\vb{\lambda}\in\C$，使
    \begin{Equation}[特征值与特征向量]
        \vb{A}\vb{v}=\lambda\vb{v}
    \end{Equation}
    则称$\lambda$是特征值（Eigenvalue），而称$\vb{v}$为与之相关的特征向量（Eigenvector）。
\end{BoxDefinition}

简单来说，特征值问题（Eigenvalue Problem）就是求解$\vb{A}\vb{v}=\lambda\vb{v}$，换言之，对于矩阵$\vb{A}$，我们发现其对某些向量的作用$\vb{A}\vb{v}$等价于对向量的缩放$\lambda\vb{v}$，这就是特征向量和特征值。需要注意的是，即便对于实矩阵$\vb{A}\in\R^{n\times n}$，其特征值$\lambda\in\C$和特征向量$\vb{v}\in\C^n$也是位于复数域的。

显然，如果$\vb{v}$是特征向量，那么$\alpha\vb{v}$也是特征向量，对于任意的$\alpha\in\C$。换言之，特征向量的模长不重要！因此，我们一般规定$\norm{\vb{v}}_2=1$，当然，这意味着$\vb{v}\in\C^n$还有复相位的自由。

\subsection{特征空间}

\begin{BoxDefinition}[特征空间]
    对于$\vb{A}\in\R^{n\times n}$，定义$\vb{A}$的$\lambda$特征空间（Eigenspace）为
    \begin{Equation}[特征空间]
        \mathcal{V}_\lambda=\qty{\vb{v}\in\C^n\mid\vb{A}\vb{v}=\lambda\vb{v}}
    \end{Equation}
\end{BoxDefinition}

矩阵$\vb{A}$的每个特征值$\lambda$都具有自己的特征空间$\mathcal{V}_\lambda$，包含特征值$\lambda$对应的全体特征向量。这很容易理解，因为特征向量可以任意缩放，若$\lambda$只对应一个$\vb{v}$，那特征空间就是$\vb{v}$所在直线。

\begin{BoxProperty}[特征空间和零空间]
    矩阵$\vb{A}\in\R^{n\times n}$的特征空间$\mathcal{V}_\lambda$可以表示为$\vb{A}-\lambda\vb{I}$的零空间
    \begin{Equation}[特征空间和零空间]
        \mathcal{V}_\lambda=\nullspace(\vb{A}-\lambda\vb{I})
    \end{Equation}
\end{BoxProperty}

这是因为求解$\vb{A}\vb{v}=\lambda\vb{v}$相当于是求解$(\vb{A}-\lambda\vb{I})\vb{v}=\vb{0}$，后者即$\vb{A}-\lambda\vb{I}$的零空间。

\begin{BoxProperty}[特征空间的不变性]
    矩阵$\vb{A}\in\R^{n\times n}$的特征空间$\mathcal{V}_\lambda$对于$\vb{A}$不变子空间
    \begin{Equation}[特征空间的不变性]
        \vb{A}\mathcal{V}_{\lambda}\subseteq\mathcal{V}_{\lambda}
    \end{Equation}
\end{BoxProperty}

这是因为$\mathcal{V}_{\lambda}$中的元素满足$\vb{A}\vb{v}=\lambda\vb{v}$，对于$\vb{A}$的变换都是缩放，因此仍位于子空间内。

\subsection{特征多项式}

\begin{BoxProperty}[特征值的数量]
    任意矩阵$\vb{A}\in\R^{n\times n}$都具有$n$个特征值。
\end{BoxProperty}

为什么$n\times n$的矩阵一定会有$n$个特征值？我们知道$\vb{A}\vb{v}=\lambda\vb{v}$等价于$(\vb{A}-\lambda\vb{I})\vb{v}=\vb{0}$，而齐次方程存在非零解$\vb{v}\neq 0$的充要条件是系数矩阵$\vb{A}-\lambda\vb{I}$奇异，即$\vb{A}-\lambda\vb{I}$的行列式为零
\begin{Equation}[特征多项式1]
    \det(\vb{A}-\lambda\vb{I})=0
\end{Equation}

我们不妨将这个表达式定义为一个关于$\lambda$的多项式函数（这里$\lambda$就是普通的变量）。

\begin{BoxDefinition}[特征多项式]
    矩阵$\vb{A}\in\R^{n\times n}$的特征多项式（Characteristic Polyomial）定义为
    \begin{Equation}[特征多项式]
        p(\lambda)=\det(\vb{A}-\lambda\vb{I})
    \end{Equation}
\end{BoxDefinition}

$p(\lambda)$是一个$n$阶的多项式，根据行列式的定义

\begin{BoxProperty}[特征多项式的阶数]
    矩阵$\vb{A}\in\R^{n\times n}$的特征多项式$p(\lambda)$是一个$n$阶多项式
    \begin{Equation}[特征多项式的阶数]
        p(\lambda)=\alpha_0+\alpha_1\lambda+\cdots+\alpha_n\lambda^n
    \end{Equation}
\end{BoxProperty}

$p(\lambda)$作为$n$阶行列式有$n$个根

\begin{BoxProperty}[特征多项式的因式分解]
    矩阵$\vb{A}\in\R^{n\times n}$的特征多项式$p(\lambda)$可以写成因式分解的形式
    \begin{Equation}[特征多项式的因式分解]
        p(\lambda)=\Prod[i=1][n](\lambda_i-\lambda)
    \end{Equation}
\end{BoxProperty}

因此，特征值的本质是特征多项式$p(\lambda)=0$的解，既然$n$阶多项式有$n$个解，那么$n\times n$的矩阵也就自然有$n$个特征值了。作为代数方程的根，我们可以预期一些特征值的性质
\begin{itemize}
    \item 特征值必然有$n$个，但可能存在重复的，这来自方程的重根。
    \item 特征值如果是复数，那一定有一对共轭的复特征值。
\end{itemize}

我们用谱来描述特征值构成的集合

\begin{BoxDefinition}[谱]
    矩阵$\vb{A}\in\R^{n\times n}$的全体特征值构成的集合称为$\vb{A}$的谱（Spectrum）
    \begin{Equation}[谱]
        \sigma(\vb{A})=\qty{\lambda\in\C\mid\det(\vb{A}-\lambda\vb{I})=0}
    \end{Equation}
\end{BoxDefinition}

\subsection{特征值的重数}

我们通常记$\lambda_1,\cdots,\lambda_n$构成的集合为$\qty{\lambda_1,\cdots,\lambda_k}$，由于特征值可能重复，有$k\leq n$成立。

\begin{BoxDefinition}[特征值的代数重数]
    定义$\lambda_i$在特征方程中的重数，为代数重数（Algebraic Multiplicity），记作$\mu_i$。
\end{BoxDefinition}

\begin{BoxDefinition}[特征值的几何重数]
    定义$\lambda_i$对应特征空间的维度，为几何重数（Geometric Multiplicity），记作$\gamma_i$
\end{BoxDefinition}

现在的问题是，代数重数$\mu_i$和几何重数$\gamma_i$是否存在某种约束关系？有趣的是，尽管两个定义来自完全不同的地方，但却存在$\mu_i\geq\gamma_i$的约束，即代数重数总是大于等于几何重数！这也就是说，如果特征值$\lambda_i$是一个$2$重根，那么特征值$\lambda_i$最多对应$2$个线性无关的特征向量。

\begin{BoxProperty}[代数重数和几何重数]
    特征值$\lambda_i$的代数重数$\mu_i$总是大于等于几何重数$\gamma_i$
    \begin{Equation}[代数重数和几何重数]
        \mu_i\geq\gamma_i
    \end{Equation}
\end{BoxProperty}

\begin{Proof}[\cref{ppt:代数重数和几何重数}]
    设矩阵$\vb{A}\in\R^{n\times n}$，设$\bar{\lambda}$是$\vb{A}$的特征值，记其几何重数为$r$
    \begin{Equation}[代数重数和几何重数1]
        r=\dim\nullspace(\vb{A}-\bar{\lambda}\vb{I})
    \end{Equation}

    现在取$n$个单位正交向量
    \begin{Equation}[代数重数和几何重数2]
        \vb{q}_1,\cdots,\vb{q}_r\in\nullspace(\vb{A}-\bar{\lambda}\vb{I})\qquad
        \vb{q}_{r+1},\cdots,\vb{q}_{n}\in\nullspace(\vb{A}-\bar{\lambda}\vb{I})^{\perp}
    \end{Equation}

    构造正交向量
    \begin{Gather}[代数重数和几何重数]
        \vb{Q}_1=\mqty[\vb{q}_1&\cdots&\vb{q}_r]\in\C^{n\times r}\id{3}\\
        \vb{Q}_2=\mqty[\vb{q}_{r+1}&\cdots&\vb{q}_n]\in\C^{n\times(n-r)}\id{4}
    \end{Gather}

    以及
    \begin{Equation}[代数重数和几何重数5]
        \vb{Q}=\mqty[\vb{Q}_1&\vb{Q}_2]
    \end{Equation}

    现在让我们计算$\vb{Q}^H\vb{A}\vb{Q}$
    \begin{Equation}[代数重数和几何重数6]
        \vb{Q}^H\vb{A}\vb{Q}=
        \mqty[
            \vb{Q}_1^H\\
            \vb{Q}_2^H\\
        ]
        \vb{A}
        \mqty[
            \vb{Q}_1&\vb{Q}_2
        ]=
        \mqty[
            \vb{Q}_1^h\vb{A}\vb{Q}_1 & \vb{Q}_1^H\vb{A}\vb{Q}_2 \\
            \vb{Q}_2^h\vb{A}\vb{Q}_1 & \vb{Q}_2^H\vb{A}\vb{Q}_2 \\
        ]
    \end{Equation}

    注意到$\vb{A}\vb{Q}_1=\bar{\lambda}\vb{Q}_1$
    \begin{Equation}[代数重数和几何重数7]
        \vb{Q}^H\vb{A}\vb{Q}=\mqty[
            \bar{\lambda}\vb{Q}_1^H\vb{Q}_1 & \vb{Q}_1^H\vb{A}\vb{Q}_2 \\
            \bar{\lambda}\vb{Q}_2^H\vb{Q}_1 & \vb{Q}_2^H\vb{A}\vb{Q}_2 \\
        ]
    \end{Equation}

    根据酉矩阵的性质，有$\vb{Q}_1^H\vb{Q}_1=\vb{I}$和$\vb{Q}_2^H\vb{Q}_1=\vb{0}$成立
    \begin{Equation}[代数重数和几何重数8]
        \vb{Q}^H\vb{A}\vb{Q}=\mqty[
            \bar{\lambda}\vb{I} & \vb{Q}_1^H\vb{A}\vb{Q}_2 \\
            \vb{0} & \vb{Q}_2^H\vb{A}\vb{Q}_2 \\
        ]
    \end{Equation}

    写出特征多项式$p(\lambda)$，注意到$\vb{Q}$不改变行列式的结果
    \begin{Equation}[代数重数和几何重数9]
        p(\lambda)=\det(\vb{A}-\lambda\vb{I})=\det(\vb{Q}^H(\vb{A}-\lambda\vb{I})\vb{Q})=\det(\vb{Q}^H\vb{A}\vb{Q}-\lambda\vb{I})
    \end{Equation}

    即有
    \begin{Equation}[代数重数和几何重数10]
        p(\lambda)=\det\mqty[
            (\bar{\lambda}-\lambda)\vb{I} & \vb{Q}_1^H\vb{A}\vb{Q}_2 \\
            \vb{0} & \vb{Q}_2^H\vb{A}\vb{Q}_2-\lambda\vb{I} \\
        ]
    \end{Equation}

    这是一个分块上三角矩阵
    \begin{Equation}[代数重数和几何重数11]
        p(\lambda)=(\bar{\lambda}-\lambda)^r\det(\vb{A}-\lambda\vb{I})
    \end{Equation}

    我们设$\bar{\lambda}$的几何重数是$r$，这里$p(\lambda)$中$(\bar{\lambda}-\lambda)^r$已经提供了$r$个$\bar{\lambda}$的重根，而其后未因式分解的$\det(\vb{A}-\lambda\vb{I})$有可能包含更多$\bar{\lambda}$的重根，这就证明了代数重数大于等于几何重数！
\end{Proof}

该证明是特征值问题的核心，具有重要意义！

\subsection{特征值的性质}

\begin{BoxProperty}[特征值和行列式]
    特征值之积，是矩阵的行列式
    \begin{Equation}[特征值和行列式]
        \det(\vb{A})=\Prod[i=1][n]\lambda_i
    \end{Equation}
\end{BoxProperty}

\begin{BoxProperty}[特征值和迹]
    特征值之和，是矩阵的迹
    \begin{Equation}[特征值和迹]
        \tr(\vb{A})=\Sum[i=1][n]\lambda_i
    \end{Equation}
\end{BoxProperty}

\begin{BoxProperty}[矩阵幂的特征值]
    矩阵$\vb{A}^k$的特征值和特征向量分别是$\lambda_1^k,\cdots,\lambda_n^k$和$\vb{v}_1^k,\cdots,\vb{v}^{n}$。
\end{BoxProperty}

\begin{BoxProperty}[矩阵的零特征值]
    矩阵$\vb{A}\in\R^{n\times n}$不满秩$\rank(\vb{A})<n$，当且仅当$0$是$\vb{A}$的特征值。
\end{BoxProperty}

通常所说的特征向量是满足$\vb{A}\vb{v}=\lambda\vb{v}$的右特征向量，也有满足$\vb{u}^T\vb{A}=\lambda\vb{u}^T$的左特征向量。不过，为此引入一个新的概念并不值得。因为$\vb{A}$的左特征向量其实就是$\vb{A}^T$的右特征向量。这也引出了另外一个论断：$\vb{A}$和$\vb{A}^T$具有相同的特征值，但一般来说对应的特征向量是不同的。
